% ── §7e Methods: Biological Interpretation of Learned Rules ──
\section{Biological Interpretation of Learned Rules}
\label{sec:methods-biological}

\subsection{Hydrophobic Core Packing}
The dominant rules discovered by the Level-1 universal circuit involve pairs
of hydrophobic or aromatic residues at short separations.  Specifically,
essential prime implicants include the conjunction
$\text{hydrophobic} \times \text{aromatic}$ at separation bins $[6, 11)$ and
$[11, 21)$, corresponding to aliphatic--aromatic packing interactions that
stabilise protein cores.  At the finer Level-2 resolution, specific residue
vocabularies emerge: $\{M, F\} \times \{M, F, Y, W\}$ and
$\{I, V\} \times \{I, V, F, W\}$ are among the highest-ranked ON cells.
These correspond precisely to the Met--Phe, Ile--Val, and Ile--Phe contacts
that dominate the buried cores of globular proteins.

\subsection{Salt Bridges and Electrostatic Interactions}
The circuit also discovers acidic--basic residue pairings at short range,
consistent with salt-bridge formation between glutamate/aspartate and
lysine/arginine side chains.  These rules appear primarily at separation bins
$[6, 11)$ and $[11, 21)$, matching the expected geometric range for
side-chain electrostatic interactions.

\subsection{Sulfur-Mediated Contacts}
Cysteine (the sole sulfur-group residue) appears in several high-tension
rules paired with hydrophobic and aromatic partners.  While disulfide bonds
require oxidative conditions and are therefore context-dependent, the
enrichment of cysteine-containing rules suggests that thiol-mediated contacts
contribute to the structural signal captured by the Karnaugh-map framework.

\subsection{Conservation and Structural Constraint}
The complete circuit's conservation dimension reveals a clear hierarchy:
pairs in which both columns are near-conserved (entropy between 0.001 and
0.05) show the strongest contact enrichment ($1.35\times$), while absolutely
conserved positions (entropy $\approx 0$) show depletion ($0.69\times$).
This non-monotonicity is biologically interpretable: absolutely invariant
positions are frequently catalytic residues or post-translational modification
sites whose conservation reflects functional constraint rather than structural
packing.  Near-conserved positions, by contrast, are sufficiently constrained
to maintain structural integrity while retaining enough variability to
generate detectable covariation signal.

\subsection{Connection to Protein Design}
The fact that Quine--McCluskey-minimised circuits discover known structural
principles from sequence data alone has implications for protein engineering.
If the learned SOP expressions capture genuine design rules---for example,
that a hydrophobic residue at position $i$ requires a hydrophobic or aromatic
partner at position $j$ within a certain sequence separation---then these
rules could serve as design constraints in inverse protein folding algorithms.
The Boolean form of these constraints makes them directly implementable as
logic circuits or constraint-satisfaction problems, offering a novel route
from coevolutionary analysis to automated protein design.
